
\documentclass{article}
\usepackage{colortbl}
\usepackage{makecell}
\usepackage{multirow}
\usepackage{supertabular}

\begin{document}

\newcounter{utterance}

\centering \large Interaction Transcript for game `cladder', experiment `full\_v1.5\_default', episode 9980 with Qwen3.5{-}27B.
\vspace{24pt}

{ \footnotesize  \setcounter{utterance}{1}
\setlength{\tabcolsep}{0pt}
\begin{supertabular}{c@{$\;$}|p{.15\linewidth}@{}p{.15\linewidth}p{.15\linewidth}p{.15\linewidth}p{.15\linewidth}p{.15\linewidth}}
    \# & \multicolumn{2}{c}{Player} && \multicolumn{2}{c}{Game Master} \\
    \hline

    \theutterance \stepcounter{utterance}  
    & & & \multicolumn{4}{p{0.6\linewidth}}{
        \cellcolor[rgb]{0.9,0.9,0.9}{
            \makecell[{{p{\linewidth}}}]{
                \texttt{\tiny{[P1$\langle$GM]}}
                \texttt{You are an expert in causal inference. The following question is not a typical commonsense query, but rather a meticulously designed question created by a professor specializing in causal inference, intended to assess the students' mastery of the course content.} \\
\texttt{Imagine a self{-}contained, hypothetical world with only the following conditions, and without any unmentioned factors or causal relationships: Hwax has a direct effect on pexu and rukz. Pexu has a direct effect on rukz. We know that hwax causes pexu. hwax and pexu causes rukz. We observed an individual is not hwax. Would an individual is not rukz if not pexu instead of pexu?} \\
\texttt{Start your answer with 'Yes' or 'No', followed by additional reasoning or evidence to support your explanation.} \\
            }
        }
    }
     \\ \\

    \theutterance \stepcounter{utterance}  
    & \multicolumn{4}{p{0.6\linewidth}}{
        \cellcolor[rgb]{0.561,0.737,0.561}{
            \makecell[{{p{\linewidth}}}]{
                \texttt{\tiny{[P1$\rangle$GM]}}
                \texttt{No} \\
\\ 
\texttt{**Reasoning:**} \\
\\ 
\texttt{To determine the answer, we must reconstruct the causal structure and the specific counterfactual scenario based on the provided conditions.} \\
\\ 
\texttt{1.  **Causal Structure**:} \\
\texttt{    *   \$Hwax \textbackslash{}to Pexu\$} \\
\texttt{    *   \$Hwax \textbackslash{}to Rukz\$} \\
\texttt{    *   \$Pexu \textbackslash{}to Rukz\$} \\
\texttt{    *   There are no other factors or relationships.} \\
\\ 
\texttt{2.  **Observed Facts**:} \\
\texttt{    *   The individual is **not** \$Hwax\$ (\$Hwax = 0\$).} \\
\texttt{    *   The individual **is** \$Pexu\$ (\$Pexu = 1\$).} \\
\texttt{    *   The individual **is** \$Rukz\$ (\$Rukz = 1\$). (Implied by the question asking "if not pexu instead of pexu" in the context of the current state, and the fact that \$Pexu\$ causes \$Rukz\$).} \\
\\ 
\texttt{3.  **The Counterfactual Question**:} \\
\texttt{    *   "Would an individual be not \$Rukz\$ if not \$Pexu\$ instead of \$Pexu\$?"} \\
\texttt{    *   Formally: Given \$Hwax = 0\$ and \$Pexu = 1\$, what is the value of \$Rukz\$ in the counterfactual world where \$Pexu = 0\$? (\$Rukz\_\{Pexu=0\}\$).} \\
\\ 
\texttt{4.  **Analysis of the Counterfactual World**:} \\
\texttt{    *   In this hypothetical world, we intervene to set \$Pexu = 0\$.} \\
\texttt{    *   The value of \$Hwax\$ remains fixed at \$0\$ (as it is an observed fact about the individual, not the variable being intervened upon).} \\
\texttt{    *   We must determine the value of \$Rukz\$ given \$Hwax = 0\$ and \$Pexu = 0\$.} \\
\texttt{    *   \$Rukz\$ is caused by two parents: \$Hwax\$ and \$Pexu\$.} \\
\texttt{    *   Since \$Hwax = 0\$, the causal path from \$Hwax\$ to \$Rukz\$ contributes nothing (assuming a standard causal model where the absence of a cause implies the absence of its effect, or at least that the cause is not active).} \\
\texttt{    *   Since we intervened to set \$Pexu = 0\$, the causal path from \$Pexu\$ to \$Rukz\$ is also inactive.} \\
\texttt{    *   With **no active causes** (\$Hwax=0\$ and \$Pexu=0\$) and **no unmentioned factors** (as per the problem statement), \$Rukz\$ cannot be generated.} \\
\\ 
\texttt{5.  **Conclusion**:} \\
            }
        }
    }
    & & \\ \\

    \theutterance \stepcounter{utterance}  
    & & & \multicolumn{2}{p{0.3\linewidth}}{
        \cellcolor[rgb]{0.9,0.9,0.9}{
            \makecell[{{p{\linewidth}}}]{
                \texttt{\tiny{[GM$|$GM]}}
                \texttt{No} \\
            }
        }
    }
    & & \\ \\

    \theutterance \stepcounter{utterance}  
    & & & \multicolumn{2}{p{0.3\linewidth}}{
        \cellcolor[rgb]{0.9,0.9,0.9}{
            \makecell[{{p{\linewidth}}}]{
                \texttt{\tiny{[GM$|$GM]}}
                \texttt{yes} \\
            }
        }
    }
    & & \\ \\

    \theutterance \stepcounter{utterance}  
    & & & \multicolumn{2}{p{0.3\linewidth}}{
        \cellcolor[rgb]{0.9,0.9,0.9}{
            \makecell[{{p{\linewidth}}}]{
                \texttt{\tiny{[GM$|$GM]}}
                \texttt{game\_result = LOSE} \\
            }
        }
    }
    & & \\ \\

\end{supertabular}
}

\end{document}
